%% ---------------------------------------------------------------------------
%%  sak-guard -- TechRxiv preprint (IEEEtran, conference option)
%%  Compile:  pdflatex paper && bibtex paper && pdflatex paper && pdflatex paper
%%  or:       latexmk -pdf paper.tex
%%
%%  References are embedded via filecontents* so this is a single-file build.
%% ---------------------------------------------------------------------------
\begin{filecontents*}{references.bib}
%% ===== Cluster D: the author's prior methodological foundation =====
@article{firmansyah2022efficacy,
  title   = {The influence of efficacy, credibility, and normative pressure to
             {M}-banking adoption level in {Indonesia}},
  author  = {Firmansyah, Irfan Achmad and Yasirandi, Rahmat and Utomo, Rio Guntur},
  journal = {Procedia Computer Science},
  volume  = {197},
  pages   = {51--60},
  year    = {2022},
  note    = {Sixth Information Systems International Conference ({ISICO} 2021)},
  doi     = {10.1016/j.procs.2021.12.117}
}

@inproceedings{firmansyah2023forensic,
  title     = {Forensic Analysis Model on Fintech E-Wallet Using Digital Evidence Generalization},
  author    = {Firmansyah, Irfan Achmad and Cahyani, Niken Dwi Wahyu and Afianti, Farah},
  booktitle = {{IEEE} Conference Proceedings},
  year      = {2023},
  address   = {Bandung, Indonesia},
  note      = {Faculty of Computing, Telkom University. DOI to be confirmed by the author}
}

%% ===== Cluster A: AI-agent security and tool-use vulnerabilities =====
@inproceedings{zhan2024injecagent,
  title     = {{InjecAgent}: Benchmarking Indirect Prompt Injections in
               Tool-Integrated Large Language Model Agents},
  author    = {Zhan, Qiusi and Liang, Zhixiang and Ying, Zifan and Kang, Daniel},
  booktitle = {Findings of the Association for Computational Linguistics: {ACL} 2024},
  year      = {2024},
  doi       = {10.18653/v1/2024.findings-acl.624}
}

@inproceedings{jia2025taskshield,
  title     = {The Task Shield: Enforcing Task Alignment to Defend Against Indirect
               Prompt Injection in {LLM} Agents},
  author    = {Jia, Feiran and Wu, Tong and Qin, Xin and Squicciarini, Anna},
  booktitle = {Proceedings of the 63rd Annual Meeting of the Association for
               Computational Linguistics ({ACL} 2025)},
  year      = {2025},
  doi       = {10.18653/v1/2025.acl-long.1435}
}

@inproceedings{greshake2023notwhat,
  title     = {Not What You've Signed Up For: Compromising Real-World
               {LLM}-Integrated Applications with Indirect Prompt Injection},
  author    = {Greshake, Kai and Abdelnabi, Sahar and Mishra, Shailesh and
               Endres, Christoph and Holz, Thorsten and Fritz, Mario},
  booktitle = {Proceedings of the 16th {ACM} Workshop on Artificial Intelligence
               and Security ({AISec} '23)},
  year      = {2023},
  doi       = {10.1145/3605764.3623985}
}

@inproceedings{cao2026survey,
  title     = {A Systematic Survey of Prompt Injection Attacks and Defenses in
               {LLM}-Integrated Applications and Agentic Systems},
  author    = {Cao, Xin and Bao, Yiping},
  booktitle = {2026 12th International Conference on Big Data and Information
               Analytics ({BigDIA})},
  year      = {2026},
  doi       = {10.1109/bigdia71496.2026.11684636}
}

@misc{zou2023universal,
  title         = {Universal and Transferable Adversarial Attacks on Aligned
                   Language Models},
  author        = {Zou, Andy and Wang, Zifan and Kolter, J. Zico and Fredrikson, Matt},
  year          = {2023},
  eprint        = {2307.15043},
  archivePrefix = {arXiv},
  primaryClass  = {cs.CL},
  doi           = {10.48550/arXiv.2307.15043}
}

@misc{chen2024struq,
  title         = {{StruQ}: Defending Against Prompt Injection with Structured Queries},
  author        = {Chen, Sizhe and Piet, Julien and Sitawarin, Chawin and Wagner, David},
  year          = {2024},
  eprint        = {2402.06363},
  archivePrefix = {arXiv},
  primaryClass  = {cs.CR},
  doi           = {10.48550/arXiv.2402.06363}
}

@article{gandhi2026vibecoding,
  title   = {Survey on Systemic Security Risks in Vibe Coding and Autonomous
             Development Workflows},
  author  = {Gandhi, Deep},
  journal = {TechRxiv},
  year    = {2026},
  doi     = {10.36227/techrxiv.176800890.09196406/v1}
}

@misc{owasp2024llm,
  title        = {{OWASP} Top 10 for Large Language Model Applications},
  author       = {{OWASP Foundation}},
  year         = {2024},
  howpublished = {\url{https://owasp.org/www-project-top-10-for-large-language-model-applications/}},
  note         = {Accessed 2026}
}

%% ===== Cluster B: static analysis and taint tracking for dynamic languages =====
@article{hou2024llmse,
  title   = {Large Language Models for Software Engineering: A Systematic
             Literature Review},
  author  = {Hou, Xinyi and Zhao, Yanjie and Liu, Yue and Yang, Zhou and
             Wang, Kailong and Li, Li and Luo, Xiapu and Lo, David and
             Grundy, John and Wang, Haoyu},
  journal = {{ACM} Transactions on Software Engineering and Methodology},
  volume  = {33},
  number  = {8},
  year    = {2024},
  doi     = {10.1145/3695988}
}

@inproceedings{yang2024falsepositives,
  title     = {Reducing False Positives of Static Bug Detectors Through Code
               Representation Learning},
  author    = {Yang, Yixin and Wen, Ming and Gao, Xiang and Zhang, Yuting and
               Jin, Yidan and Wang, Haoyu and Chen, Zhenyu},
  booktitle = {2024 {IEEE} International Conference on Software Analysis,
               Evolution and Reengineering ({SANER})},
  year      = {2024},
  doi       = {10.1109/saner60148.2024.00075}
}

@article{borlakova2026comparative,
  title   = {Comparative evaluation of static analysis of program code using
             {CodeQL}, {Semgrep} and {SonarQube}},
  author  = {Borlakova, Mariam A. and Sitnikov, Sergey Yu. and Vilkov, Vladislav E.},
  journal = {Soft Measurements and Computing},
  year    = {2026},
  doi     = {10.36871/2618-9976.2026.08.002}
}

@inproceedings{liu2025bidirectional,
  title     = {Bi-directional Taint Flow Analysis: A High-precision Static
               Detection Approach for Java Deserialization Vulnerabilities},
  author    = {Liu, Haihua and Lu, Yanrong},
  booktitle = {2025 4th International Symposium on Computer Applications and
               Information Technology ({ISCAIT})},
  year      = {2025},
  doi       = {10.1109/iscait64916.2025.11010573}
}

@article{shimchik2019taint,
  title   = {Vulnerabilities Detection via Static Taint Analysis},
  author  = {Shimchik, N. V. and Ignatyev, V. N.},
  journal = {Proceedings of the Institute for System Programming of the {RAS}},
  volume  = {31},
  number  = {3},
  year    = {2019},
  doi     = {10.15514/ispras-2019-31(3)-14}
}

@misc{bandit2024,
  title        = {{Bandit}: A security linter for {Python} source code},
  author       = {{PyCQA}},
  year         = {2024},
  howpublished = {\url{https://github.com/PyCQA/bandit}},
  note         = {Accessed 2026}
}

%% ===== Cluster C: blockchain client-side and smart-contract security =====
@inproceedings{zhao2025soleker,
  title     = {Soleker: Uncovering Vulnerabilities in {Solana} Smart Contracts},
  author    = {Zhao, Kunsong and Tian, Yunpeng and Ma, Zuchao and Luo, Xiapu},
  booktitle = {2025 40th {IEEE}/{ACM} International Conference on Automated
               Software Engineering ({ASE})},
  year      = {2025},
  doi       = {10.1109/ase63991.2025.00209}
}

@inproceedings{yang2024sokai,
  title     = {{SoK}: {AI}-Powered Security Analysis of Smart Contract},
  author    = {Yang, Shuo and Niu, Jianyu and Zhang, Yinqian},
  booktitle = {2024 International Conference on Meta Computing ({ICMC})},
  year      = {2024},
  doi       = {10.1109/icmc60390.2024.00017}
}

@article{wijaya2025ethereum,
  title   = {Ethereum vs {Solana}: A Comparative Study of Blockchain Architecture
             on Performance, Security, and Ecosystem Development},
  author  = {Wijaya, Michael and Melvyn, Franscelino and Setiawan, Reina and
             Rumagit, Reinert Yosua},
  journal = {Procedia Computer Science},
  year    = {2025},
  doi     = {10.1016/j.procs.2025.08.273}
}

@article{khan2024survey,
  title   = {A Survey of Vulnerability Detection Techniques by Smart Contract Tools},
  author  = {Khan, Zulfiqar Ali and Namin, Akbar Siami},
  journal = {{IEEE} Access},
  volume  = {12},
  year    = {2024},
  doi     = {10.1109/access.2024.3401623}
}

@inproceedings{wang2023evaluation,
  title     = {Evaluation of Contemporary Smart Contract Analysis Tools},
  author    = {Wang, Baocheng and Chen, Shiping and Wang, Qin},
  booktitle = {Proceedings of the 18th International Conference on Evaluation of
               Novel Approaches to Software Engineering ({ENASE})},
  year      = {2023},
  doi       = {10.5220/0011769600003464}
}

@article{venkatraman2026defi,
  title   = {Securing {DeFi}: a comprehensive review of machine learning approaches
             for detecting smart contract vulnerabilities and threats},
  author  = {Venkatraman, Dhivyalakshmi and Kuppusamy, Manikandan},
  journal = {International Journal of Informatics and Communication Technology},
  volume  = {15},
  number  = {1},
  year    = {2026},
  doi     = {10.11591/ijict.v15i1.pp438-446}
}

@misc{yakovenko2018solana,
  title        = {{Solana}: A new architecture for a high performance blockchain --
                  {Proof} of {History}, {Tower} {BFT} and {Turbine}},
  author       = {Yakovenko, Anatoly},
  year         = {2018},
  howpublished = {Solana Labs whitepaper},
  note         = {Accessed 2026}
}

@inproceedings{tantikul2020exploring,
  title     = {Exploring Vulnerabilities in Solidity Smart Contract},
  author    = {Tantikul, Phitchayaphong and Ngamsuriyaroj, Sudsanguan},
  booktitle = {Proceedings of the 6th International Conference on Information
               Systems Security and Privacy ({ICISSP})},
  year      = {2020},
  doi       = {10.5220/0008909803170324}
}

@misc{tob2024solanalints,
  title        = {{solana-lints}: {Rust} lints for {Solana} programs},
  author       = {{Trail of Bits}},
  year         = {2024},
  howpublished = {\url{https://github.com/crytic/solana-lints}},
  note         = {Accessed 2026}
}
\end{filecontents*}

\documentclass[conference]{IEEEtran}
\IEEEoverridecommandlockouts

\usepackage{cite}
\usepackage{amsmath,amssymb,amsfonts}
\usepackage{algorithmic}
\usepackage{graphicx}
\usepackage{textcomp}
\usepackage{xcolor}
\usepackage{booktabs}
\usepackage{url}
\usepackage[hidelinks]{hyperref}

\def\BibTeX{{\rm B\kern-.05em{\sc i\kern-.025em b}\kern-.08em
    T\kern-.1667em\lower.7ex\hbox{E}\kern-.125emX}}

\begin{document}

\title{AST-Driven Static Analysis and Taint Tracking for Autonomous Agent
Execution Layers in Decentralized Systems:\\
The \texttt{sak-guard} Pre-Commit Guardrail for \mbox{Solana} Agent Kit}

\author{\IEEEauthorblockN{Irfan Achmad Firmansyah}
\IEEEauthorblockA{\textit{Pantronux Research}\\
Indonesia\\
GitHub: Pantronux-Research/sak-guard\\
}
}

\maketitle

\begin{abstract}
Large language model (LLM) driven agents increasingly hold cryptographic key
material and submit transactions autonomously, yet the security tooling for
decentralized systems remains concentrated on-chain: Rust and Anchor program
analysers inspect what executes inside the virtual machine, and runtime
transaction firewalls inspect what is about to be submitted. Neither observes
the Python client and agent layer where keypairs are sourced, where transaction
construction is silently misconfigured, and where unsanitised model output
reaches a remote procedure call. We present \texttt{sak-guard}, a static
pre-commit guardrail that operates on that layer, and a nine-pattern taxonomy
(\mbox{PS-SOL-1} to \mbox{PS-SOL-3}) that formalises its defect classes. We give
an AST-gated detection formulation that eliminates a class of lexical
false positives which a naive regular-expression bank cannot avoid, and a
single-file taint propagation rule that follows model-derived data through
\texttt{json.loads} subscripts into transaction sinks. Evaluated against 32
real tool implementations (5{,}670 lines) drawn from the upstream
\texttt{sendaifun/solana-agent-kit-py} repository rather than synthetic
fixtures, the engine reaches 73\% precision (8 of 11 findings independently
judged correct), with the two pattern families that have no equivalent in
existing Solana tooling producing zero false positives. We report the
precision trajectory (41\% $\rightarrow$ 57\% $\rightarrow$ 73\%) and every
defect the evaluation surfaced in our own tool, including three remaining
false positives concentrated in a single temporal heuristic, which we retain
at \textit{medium} severity rather than suppress.
\end{abstract}

\begin{IEEEkeywords}
static analysis, taint tracking, abstract syntax tree, autonomous agents,
large language models, {Solana}, software security, pre-commit guardrails
\end{IEEEkeywords}

%% ===========================================================================
\section{Introduction}
%% ===========================================================================

The transition from LLMs as text generators to LLMs as \emph{actors} is the
defining architectural shift in contemporary software engineering
\cite{hou2024llmse}. An agent no longer merely produces a completion; it
selects a tool, constructs its arguments, and executes it. When the domain is a
decentralized system, those tools hold the keys. A single agent-authored
function may read a keypair, assemble a transaction, and submit it, and the
entire trust boundary of the resulting system rests on whether that function
does what its natural-language description claims.

The security community has developed two mature responses, and both observe
the wrong layer. \emph{On-chain} analysis, exemplified by Rust and Anchor
linters and by empirical vulnerability studies of Solana programs
\cite{zhao2025soleker,tob2024solanalints}, examines code that runs inside the
virtual machine. It is structurally incapable of observing how a Python client
obtained a signing key. \emph{Runtime} firewalls, such as the transaction
interception plugin proposed for the Solana Agent Kit, examine a transaction at
the moment of submission. By then the vulnerable code has been merged,
reviewed, and deployed; the firewall can refuse one transaction, but it cannot
repair the pattern that generated it.

Between these two sits an unguarded layer. A tool implementation in the Solana
Agent Kit is ordinary Python. It imports \texttt{solders} or \texttt{solana-py},
it reads configuration, and it calls a model. The defect classes that matter
there are not on-chain defects and are not visible to a transaction firewall,
because they describe \emph{how source code is written} rather than what
executes. A hardcoded mnemonic is a defect the moment it is typed. A
\texttt{send\_raw\_transaction} call issued without a commitment level and
without confirmation is a defect in the source, whatever the runtime does next.

We present \texttt{sak-guard}, a static pre-commit guardrail for that layer, and
make four contributions:

\begin{enumerate}
\item A nine-pattern taxonomy of client-side agent defects organised into three
  families (\S\ref{sec:taxonomy}), spanning key-material handling,
  \texttt{solders}/\texttt{solana-py} misuse, and model-output taint.
\item An AST-gated detection formulation (\S\ref{sec:engine}) that resolves
  operator tokens against the parse tree rather than the raw line, and a
  single-file taint rule that propagates provenance through subscript and
  attribute access.
\item An empirical evaluation on 32 real upstream tool files (5{,}670 lines)
  reporting per-pattern precision rather than an aggregate
  (\S\ref{sec:eval}).
\item A deliberately unflattering account of the three false positives and the
  four defects the evaluation exposed in our own implementation, retained
  inside the published figures (\S\ref{sec:limits}).
\end{enumerate}

%% ===========================================================================
\section{Related Work and Theoretical Framing}
%% ===========================================================================

\subsection{Precedents in the Author's Prior Work}

This work did not begin from the tooling literature; it descends from two
earlier studies of system behaviour under adversarial conditions, and the
methodological lineage is direct. In \cite{firmansyah2022efficacy} we modelled
adoption of a transaction-bearing consumer system through the Technology
Acceptance Model, establishing that users' willingness to entrust value to an
automated channel is governed by \emph{perceived credibility} and
\emph{self-efficacy} as much as by functional utility. The operative finding for
the present work is that confidence in a value-bearing system is a function of
its observable trust properties, not of its capability. An autonomous agent
that can move funds is precisely such a system, and its credibility is
determined by properties a reviewer can inspect in the source.

In \cite{firmansyah2023forensic} we developed a forensic analysis model for
fintech e-wallet applications, using an RDF ontology to generalise acquired
digital evidence across User, Transaction and Merchant entities so that
investigative questions could be answered from heterogeneous application data.
That work formalised the same underlying problem this paper addresses:
\emph{heterogeneous, application-specific artefacts must be mapped onto a
common abstraction before anomalies become detectable}. There, the abstraction
was an evidence ontology; here, it is a pattern taxonomy and a dataflow
relation. The shift from post-incident forensic reconstruction to pre-commit
static prevention is the contribution this paper adds: the earlier work asked
what evidence a compromised system leaves behind, and this work asks how to
refuse the defect that would generate it.

\subsection{Agent and Tool-Use Security}

The vulnerability class that motivates agent-layer analysis is indirect prompt
injection: adversarial content in data the agent consumes redirects its
tool use \cite{greshake2023notwhat}. \textsc{InjecAgent} operationalises this
for tool-integrated agents, demonstrating that attacker-controlled inputs
reliably induce unintended tool invocation \cite{zhan2024injecagent}, and
\cite{chen2024struq} and \cite{jia2025taskshield} propose architectural and
task-alignment defences. Recent surveys consolidate the attack surface and its
mitigations \cite{cao2026survey}, and the {OWASP} Top 10 for LLM Applications
gives the practitioner taxonomy \cite{owasp2024llm}. Adversarial-input work
establishes that alignment does not survive optimisation \cite{zou2023universal},
which is the assumption under which our third pattern family operates: we do not
attempt to make the model trustworthy, we assert that its output is
untrusted \emph{at the sink}. A closely related study documents systemic
security risk in AI-assisted development workflows \cite{gandhi2026vibecoding},
and reaches the same conclusion this paper is built on: when implementation pace
is accelerated by model authorship, the review gate must be automated, because
manual review does not scale to the rate at which the defects are produced.

\subsection{Static Analysis and Taint Tracking}

Static detection in dynamic languages faces an inherent precision problem, and
this is a well-characterised field. Taint analysis is the established technique
for tracking untrusted provenance to a sensitive sink
\cite{shimchik2019taint,liu2025bidirectional}, and modern tooling such as
Bandit \cite{bandit2024} and Semgrep applies pattern and AST matching at
practitioner scale, with comparative evaluations of these engines now available
\cite{borlakova2026comparative}. The persistent difficulty is false positives:
a linter that fires on correct code is abandoned, which is a failure mode
indistinguishable from never firing at all. Dedicated work attacks this
directly through learned code representations \cite{yang2024falsepositives},
and our evaluation design (\S\ref{sec:eval}) adopts its premise that
precision must be \emph{measured} on real code rather than asserted.

We differ from these tools in scope rather than method. Bandit and Semgrep are
general-purpose; they have no notion of a \texttt{solders} commitment level, of
a blockhash's temporal validity, or of a Pubkey sink. Our contribution is a
domain taxonomy with the same static-analysis discipline applied to it, and an
evaluation that isolates which patterns generalise and which do not.

\subsection{On-Chain and Client-Side Blockchain Security}

Smart-contract vulnerability detection is mature and heavily surveyed
\cite{khan2024survey,venkatraman2026defi,yang2024sokai}, and tool evaluations
show broad but uneven coverage across the mainstream toolset
\cite{wang2023evaluation,tantikul2020exploring}. Solana's architecture, with
parallel execution, Proof of History and a distinct account model
\cite{yakovenko2018solana,wijaya2025ethereum}, requires dedicated treatment,
and recent work does exactly that at the program level \cite{zhao2025soleker}
while the ecosystem's own lints target Rust programs \cite{tob2024solanalints}.
The gap this paper fills is consistent across this literature: \emph{every}
surveyed tool operates on the contract or the program, and none addresses the
Python client that an autonomous agent uses to reach it. Our taxonomy is
therefore complementary to this body of work rather than competitive with it,
and we position \texttt{sak-guard} as the static counterpart to a runtime
transaction firewall: the firewall decides whether a given transaction may
execute, and we decide whether the code that builds it should be merged.

%% ===========================================================================
\section{Threat Model and Vulnerability Taxonomy}
\label{sec:taxonomy}
%% ===========================================================================

\subsection{Threat Model}

We consider an autonomous agent that holds a signing keypair and calls
tools to construct and submit transactions. We assume: (i) the model is
adversarially influenceable, so its output is untrusted
\cite{zou2023universal,greshake2023notwhat}; (ii) the attacker may contribute
the agent's tool implementations or their natural-language descriptions, in
which case a defect can be introduced as ordinary-looking code but not as an
obvious one; (iii) the on-chain program itself is out of scope. Under these
assumptions the adversary's goal is to obtain key material, to cause a
transaction whose semantics differ from the stated intent, or to reach an RPC
sink with attacker-directed arguments. Notably, every one of these goals can be
achieved by code that passes type checking and unit tests, which is why static
detection at the source is a necessary control rather than a redundant one.

\subsection{PS-SOL-1: Unsafe Key Material in Agent Code}

The first family covers literal key material or key material logged or
embedded by construction. We define four sub-patterns. \textsc{1a} flags a
base58 literal passed to \texttt{Keypair.from\_base58\_string}; \textsc{1b} a
seed literal passed to \texttt{from\_seed}; \textsc{1c} a 12- or 24-word
mnemonic literal; \textsc{1d} a key-named identifier interpolated into an
f-string. Formally, for a call node $c$ with callee $f(c)$ and argument set
$A(c)$:

\begin{equation}
\mathrm{leak}(c) \iff f(c) \in \mathcal{K} \;\wedge\;
\exists\, a \in A(c): \mathrm{literal}(a),
\end{equation}

where $\mathcal{K}$ is the set of key-sourcing entry points
($\texttt{from\_base58\_string}$, $\texttt{from\_seed}$, $\texttt{Keypair}$).
For \textsc{1d} a name-based predicate is required, and we note its weakness
explicitly: identifier matching is lexical, so it can fire on a correctly
handled key. The mitigation is semantic rather than lexical---if the
interpolated expression is itself a public-key accessor
(\texttt{.pubkey()}, \texttt{.public\_key()}, \texttt{.to\_string()}) the
value is intended to be public, and the pattern is suppressed. This
refinement was forced by our own evaluation (\S\ref{sec:eval}).

\subsection{PS-SOL-2: \texttt{solders}/\texttt{solana-py} Misuse}

The second family is behavioural: each sub-pattern describes a divergence
between what the code appears to guarantee and what the library actually
guarantees at the call site. \textsc{2a} flags a submission issued without an
explicit commitment level and without a confirmation call; the return value is
a signature for a transaction whose finality is unverified, so success is
reported optimistically. \textsc{2b} flags a
\texttt{simulate\_transaction} result whose \texttt{.err} field is never
inspected, meaning simulation is performed and then discarded, which is worse
than not simulating because it implies a check that does not exist. \textsc{2c}
flags a blockhash obtained at line $b$, followed by intervening work, followed
by a submission at line $s$:

\begin{equation}
\mathrm{stale}(b,s) \iff b < s \;\wedge\; \exists\, g:\;
b < \mathrm{line}(g) < s \;\wedge\; g \in \mathcal{S},
\end{equation}

where $\mathcal{S}$ is the set of slow-call vocabulary ($\texttt{send}$,
\texttt{create}, \texttt{sign}, \texttt{confirm}, \texttt{await}). A Solana
blockhash is valid for a bounded window; work performed between obtaining it
and submitting can exhaust that window and cause silent rejection. This is the
one pattern in the taxonomy whose detection is \emph{temporal} rather than
structural, and it is correspondingly the weakest; \S\ref{sec:eval} quantifies
the cost.

\subsection{PS-SOL-3: Unsanitised Model Output Reaching a Sink}

The third family operationalises the assumption that model output is
untrusted. We define a taint relation $T \subseteq \mathrm{Var}$ over a
single file. A variable is tainted, written $x \in T$, if it is assigned from
a model-invocation call:

\begin{equation}
\mathrm{source}(x) \iff x \leftarrow f(\cdot), \;
f \in \mathcal{L},
\end{equation}

with $\mathcal{L}$ the model-invocation vocabulary matched on \emph{stems}
rather than whole words (\texttt{complete}, \texttt{complet},
\texttt{invoke}, \texttt{chat}, \texttt{generate}, \texttt{prompt}). Stem
matching is required because a whole-word test for \texttt{completion} does not
match the callee \texttt{complete}. Taint then propagates through intermediate
parsing and access operations:

\begin{equation}
x \in T \;\wedge\; y \leftarrow \mathrm{peel}(x) \;\Rightarrow\; y \in T,
\end{equation}

where $\mathrm{peel}$ denotes subscript or attribute access. This rule is
essential and was validated empirically: the dominant real-world shape is not
a direct call but \texttt{json.loads(plan)["recipient"]}, which parses to a
\texttt{Subscript} node rather than a \texttt{Call}. A tracker that inspects
only \texttt{Assign(value=Call)} reports \emph{clear} on vulnerable code, a
silent false negative of exactly the kind this paper argues against. Finally, a
sink violation is reported when a tainted variable reaches a sensitive
position:

\begin{equation}
\mathrm{violation} \iff \exists\, x \in T,\; c:\;
x \in A(c) \;\wedge\; f(c) \in \mathcal{P},
\end{equation}

with $\mathcal{P} = \{\texttt{Pubkey}, \texttt{transfer},
\texttt{send\_transaction}, \texttt{send\_raw\_transaction}\}$.

%% ===========================================================================
\section{Detection Engine Architecture}
\label{sec:engine}
%% ===========================================================================

\subsection{Lexical Disambiguation by Parse-Tree Gating}

A detection bank expressed as line regular expressions is compact and fast,
and it is wrong in a specific and instructive way. Consider a rule intended to
find division by an unchecked variable:

\begin{equation}
\mathrm{flag}(\ell) \iff \exists\, m:\; \mathrm{match}(r_{\div}, \ell)
\;\wedge\; \neg\,\mathrm{guarded}(\ell).
\end{equation}

Applied to the \emph{raw text} of a line, the operator token is
indistinguishable from the same character inside a string literal. The
sequence \texttt{"https://api.mainnet-beta.solana.com"} contains a slash
immediately followed by an identifier, and matches $r_{\div}$; so do file paths
and compiled regular expressions. A URL therefore reads as a division. We
observed precisely this failure: an RPC endpoint initialisation was reported as
\emph{division by variable without zero-check}.

The resolution is not to make $r_{\div}$ cleverer, because the distinction is
not lexical. It is resolved by gating the rule on the parse tree:

\begin{equation}
\mathrm{flag}(\ell) \iff \ell \in L_{\div}(P) \;\wedge\;
\mathrm{match}(r_{\div},\ell) \;\wedge\; \neg\,\mathrm{guarded}(\ell),
\end{equation}

where $L_{\div}(P) = \{\mathrm{lineno}(n) : n \in P,\;
n = \mathrm{BinOp}(\mathrm{op} = \mathrm{Div})\}$ is precomputed from the AST
$P$ of the file. Two refinements are necessary. A division whose right operand
is a literal or an uppercase module constant
(\texttt{LAMPORTS\_PER\_SOL}) cannot be zero at runtime and is excluded;
and when $P$ does not parse the precomputed set is undefined and the rule
falls back to its former behaviour, so a syntax error degrades precision
rather than silently disabling the check. We state the general principle this
instantiates, because it recurred three times in our implementation:
\emph{a lexically-specified operator that is structurally disambiguable must be
gated on the parser, not on a better regular expression.}

\subsection{Pattern Execution with Verdicts and Error Isolation}

Every registered pattern emits an explicit verdict for every target:
\texttt{found}, \texttt{clear}, or \texttt{n/a}, where \texttt{n/a} means the
pattern \emph{cannot} apply because the file contains none of the domain's
marker tokens. This three-valued discipline exists because a report that is
silent about a pattern is indistinguishable from a pattern that never ran. A
file with no Solana or model-invocation vocabulary therefore returns
\texttt{n/a} for the taxonomy rather than \texttt{clear}, since \texttt{clear}
on an unrelated file would be a false claim of coverage. Each pattern executes
inside its own error boundary and records an exception as that pattern's
verdict, so a malformed input cannot take down the run: the engine never
imports or executes the code under audit, and malformed, empty, binary and
comment-only inputs all produce diagnostics rather than a traceback.

\subsection{Strategy Selection and Severity Honesty}

The base analysis bank covers credential literals and \texttt{eval}, and does
\emph{not} contain shell-injection patterns at all; those reside in the
security-focused strategy. The command-line interface therefore defaults to the
security-focused strategy, since a security auditor defaulting to a bank that
lacks command-injection detection reports \emph{clean} on
\texttt{subprocess.call(cmd, shell=True)}. Severity is assigned to match
confidence: the AST-gated key-material and taint patterns receive
\texttt{critical}; the structural library-misuse patterns receive
\texttt{high}; and the temporal heuristic \textsc{2c} receives \texttt{medium}.
The last assignment is deliberate and is defended by the measurement in the
following section.

%% ===========================================================================
\section{Empirical Evaluation}
\label{sec:eval}
%% ===========================================================================

\subsection{Corpus and Method}

The corpus is the upstream repository \texttt{sendaifun/solana-agent-kit-py} at
its \texttt{main} revision: 32 tool implementation files, 5{,}670 lines,
fetched unmodified. We deliberately did not author lookalike samples. Samples
shaped by the same assumptions that shaped the patterns cannot contradict them,
so a self-authored corpus cannot measure false-positive load; this follows the
methodological requirement established for learned false-positive reduction
\cite{yang2024falsepositives} and the same reasoning that motivated
ontology-driven evidence generalisation in our forensic work
\cite{firmansyah2023forensic}. The corpus is third-party code under its own
licence and is consequently not vendored into our repository; the revision is
recorded so the run is reproducible.

Each finding was then classified by hand as a true or false positive by
reading the flagged line in its function context. In accordance with
\cite{yang2024falsepositives} we report precision per pattern rather than in
aggregate, because the aggregate conceals exactly what a practitioner needs to
decide which rule to trust.

\subsection{Precision Trajectory}

Table~\ref{tab:traj} gives the three stages of the evaluation. Each stage was
driven by defects that the previous stage exposed, not by tuning toward a
target; every improvement came from repairing our own implementation.

\begin{table}[t]
\caption{Precision across three evaluation stages (32 files, 5{,}670 lines)}
\label{tab:traj}
\centering
\begin{tabular}{@{}lccc@{}}
\toprule
\textbf{Metric} & \textbf{Baseline} & \textbf{After FP} & \textbf{After base-} \\
                &                  & \textbf{fixes}    & \textbf{bank fixes} \\
\midrule
Findings                & 17 & 14 & \textbf{11} \\
Files with zero findings & 23 & 25 & \textbf{25} \\
\quad(\% of corpus)      & 72\% & 78\% & \textbf{78\%} \\
Judged true positives   & 7  & 8  & \textbf{8} \\
Judged false positives  & 10 & 6  & \textbf{3} \\
\textbf{Precision}      & \textbf{41\%} & \textbf{57\%} & \textbf{73\%} \\
\bottomrule
\end{tabular}
\end{table}

\subsection{Per-Pattern Outcomes}

Table~\ref{tab:pattern} disaggregates the final stage, which is the
substantive result.

\begin{table}[t]
\caption{Per-pattern outcomes on the real-code corpus}
\label{tab:pattern}
\centering
\begin{tabular}{@{}llcc@{}}
\toprule
\textbf{Pattern} & \textbf{Class} & \textbf{TP} & \textbf{FP} \\
\midrule
PS-SOL-1 & Key material (literal)      & 0 & \textbf{0} \\
PS-SOL-2a & Send without confirmation   & 1 & 0 \\
PS-SOL-2b & Simulation result ignored   & 0 & \textbf{0} \\
PS-SOL-2c & Stale blockhash ordering    & 3 & 2 \\
PS-SOL-3 & Model taint to sink         & 0 & \textbf{0} \\
\midrule
base:no-return    & Missing return        & 1 & 0 \\
base:division     & Unchecked divisor     & 2 & 1 \\
base:bare-except  & Bare \texttt{except}  & 1 & 0 \\
\midrule
\textbf{Total}    &                       & \textbf{8} & \textbf{3} \\
\bottomrule
\end{tabular}
\end{table}

The two families that exist in no other Solana tooling produced
\emph{zero} false positives across 5{,}670 lines, which is the central
empirical claim of this paper. Conversely, \textsc{2c} achieved 3 of 5, and
both of its errors share one cause: a blockhash fetched on the line
immediately preceding the submission is still counted as having intervening
work, because the submission call itself matches the slow-call vocabulary of
(\ref{eq:stale}). We retain this pattern at \texttt{medium} severity rather
than deleting it, because its two true positives are expensive defects---in
\texttt{deploy\_token.py} the blockhash is fetched at line 48, an account is
created at line 54, and submission occurs at line 92, which is genuinely
stale; in \texttt{use\_raydium.py} the interval spans lines 152 to 280. The
practitioner guidance follows from the measurement and is stated in the
artifact: a \texttt{medium} \textsc{2c} finding should be read as \emph{check
this ordering}, not as a confirmed defect.

\subsection{True Positives in Upstream Code}

Beyond the key-material and taint families, the engine identified genuine
misuse in maintained third-party code. In \texttt{transfer.py} line 67 a
submission is issued via \texttt{send\_raw\_transaction} with neither an
explicit commitment nor a confirmation call, so the returned signature
describes a transaction whose finality is unverified. In
\texttt{burn\_and\_close\_account.py} line 83 a function returning
\texttt{None} implicitly is flagged by the missing-return rule. In
\texttt{get\_tps.py} line 124 a transaction count is divided by a sample
period without a zero check. That these findings are correct is unsurprising;
that they were produced \emph{by a static rule authored against a taxonomy} in
a repository we did not write is the relevant evidence.

\subsection{Defects the Evaluation Exposed in Our Own Implementation}

Reporting only the detection rate would overstate coverage, and for a security
tool that overstatement is itself the vulnerability. The evaluation found
\emph{more} defects in \texttt{sak-guard} than in the corpus. Three were fixed
and pushed:

\begin{enumerate}
\item \textbf{Slash-in-literal read as division.} The rule in
  (\ref{eq:rawflag}) flagged an RPC endpoint literal. Fixed by the parse-tree
  gating of (\ref{eq:gated}).
\item \textbf{Public-key logging read as a key leak.} \textsc{1d} fired on
  \texttt{logger.info(f"\{kp.pubkey()\}")}. A public key is intended to be
  printed; the suppression described in \S\ref{sec:taxonomy} was added.
\item \textbf{Division by a module constant.} \texttt{x / LAMPORTS\_PER\_SOL}
  is fixed at import and never zero. Literal and uppercase-constant divisors
  are now excluded, which removed three of five division findings.
\end{enumerate}

Two further defects were found by the downstream release audit and are fixed
in the final artifact: a path argument that did not exist caused the tool to
exit \textbf{0} and report success on a source tree it never read, and the
command-line help advertised the program name of the internal tool from which
this was decoupled. The first is instructive beyond this project: an exit code
derived from \emph{findings present} returns success when the input silently
yields no files, so the documented usage-error contract was violated in a way
no test written against real files could detect. The general lesson is that a
detection tool's exit-code contract must be asserted for inputs that are
\emph{required to fail}, exactly as a detection rule must be tested against
inputs that are required to stay silent.

%% ===========================================================================
\section{Limitations and Threats to Validity}
\label{sec:limits}
%% ===========================================================================

We state the boundaries of the technique, since a static tool that claims more
than it delivers is rejected on inspection, and for a security artifact that
rejection is deserved.

\paragraph{Intra-procedural and intra-file scope}
The taint relation of (\ref{eq:taint})--(\ref{eq:violation}) is computed within
a single file. A value validated in another module, or a timestamped
simulation check performed inside a decorator, will be reported as a violation
even when the code is correct, and---more dangerously---a value round-tripped
through a database is not tracked at all. We therefore describe \textsc{3} as
\emph{single-file taint}, not whole-program dataflow, and it detects the naked
shape rather than the validated-but-lookalike shape. Extending to
inter-procedural analysis is the obvious next increment and would require
call-graph construction the current engine does not perform.

\paragraph{Temporal heuristics}
\textsc{2c} is order-based and, as \S\ref{sec:eval} quantifies, its two errors
and its potential false negatives share one root cause: ordering is a proxy for
elapsed time, and the absence of intervening calls does not establish that the
blockhash is fresh. A genuinely stale blockhash reused across functions may
also escape detection. We report it as a heuristic at \texttt{medium} severity
rather than as a defect class, and the measured 3-of-5 outcome is published
rather than smoothed.

\paragraph{Corpus and generality}
All quantitative results derive from one upstream repository chosen because it
is the exact target domain. The corpus contains no occurrence of some patterns
(the \texttt{n/a} and zero-count rows of Table~\ref{tab:pattern}), so those
rows establish the absence of false positives and \emph{not} a sensitivity or
recall figure. We do not report recall: no ground-truth vulnerability
labelling exists for these files, and precision measured on a real corpus plus
fixture-verified per-pattern sensitivity is what the evidence supports. A
larger multi-repository corpus with adjudicated labels is required before any
recall claim would be defensible.

\paragraph{Hand classification}
True/false judgements were made by the authors, which introduces the
possibility of bias toward our own patterns. We mitigate this by publishing
every finding with its line and file, so the classification is falsifiable by
any reader against the recorded revision---the same transparency argument that
motivated generalised, inspectable evidence in our forensic work
\cite{firmansyah2023forensic}. Independent adjudication is the correct
strengthening and is left as future work.

\paragraph{Baseline-threat scope}
The engine inspects source. It cannot detect a defect introduced at runtime,
nor a well-formed tool whose behaviour diverges from its description only under
adversarial input---by construction, that is the runtime firewall's task. This
is why we position the work as complementary rather than as a replacement, and
we consider the taxonomy's value to lie in moving detection to the earliest
moment at which the defect exists.

%% ===========================================================================
\section{Conclusion and Future Work}
%% ===========================================================================

We presented \texttt{sak-guard}, a static pre-commit guardrail for the Python
client and agent layer of decentralized systems, with a nine-pattern taxonomy,
an AST-gated detection formulation and a single-file taint rule. Evaluated on
32 real upstream tool files (5{,}670 lines), the engine reaches 73\% precision
with zero false positives in the two pattern families that have no equivalent
in existing Solana tooling, following a trajectory from 41\% to 57\% to 73\% in
which every gain came from repairing defects that evaluation exposed in our own
implementation. The paper's methodological claim is that this trajectory, and
the three remaining false positives, are more useful to a reader than a
detection rate would have been: the per-pattern split is what tells a
practitioner which rules to trust.

The work descends directly from our prior research. The adoption study
\cite{firmansyah2022efficacy} established that trust in a value-bearing
automated system is a function of observable trust properties; the forensic
model \cite{firmansyah2023forensic} established that heterogeneous artifacts
must be mapped to a common abstraction before anomalies are detectable. This
paper is the synthesis of both for the autonomous-agent setting, and it
inverts the earlier question: rather than reconstructing what a compromised
system left behind, we refuse the code that would compromise it.

Four directions follow. First, inter-procedural and cross-module taint
propagation via call-graph construction, which is the single largest source of
the limitations in \S\ref{sec:limits}. Second, a multi-repository corpus with
independently adjudicated labels, without which no recall claim is defensible.
Third, and most promising, the extension of the intent-mismatch dimension: the
engine currently treats model output as untrusted data, and the natural
generalisation is to compare an implementation against its declared
natural-language description and fail the build on divergence---a static
detection of the very class of attack that \cite{greshake2023notwhat} and
\cite{zhan2024injecagent} demonstrate at runtime. Fourth, integration of the
taxonomy as a shared rule format capable of expressing both static and runtime
invariants, so that the pre-commit and transaction-firewall layers can agree on
a single definition of a defect. The artifact, its specification and its full
verification report are public at
\url{https://github.com/Pantronux-Research/sak-guard}.

\section*{Acknowledgment}

The author thanks Telkom University for the research training in which the
methodological foundations of this work, and the prior publications cited
therein, were developed.

\bibliographystyle{IEEEtran}
\bibliography{references}

\end{document}
